% Part: first-order-logic
% Chapter: axiomatic-deduction
% Section: proof-theoretic-notions

\documentclass[../../../include/open-logic-section]{subfiles}

\begin{document}

\iftag{FOL}
      {\olfileid{fol}{axd}{ptn}}
      {\olfileid{pl}{axd}{ptn}}

\olsection{सिद्धता-उपपत्तीय संकल्पना}

\begin{explain}
जशा आपण अनेक महत्त्वाच्या चिन्हार्थविषयक संकल्पना
(\iftag{FOL}{वैधता}{सर्वतः सत्य सूत्र}, तार्किक निष्पन्नता, पूर्ततायोग्यता)
परिभाषित केल्या आहेत, तशाच आता त्यांना अनुरूप
\emph{सिद्धता-उपपत्तीय संकल्पना} परिभाषित करू. या संकल्पना !!{structure}s
मध्ये !!{sentence}s ची पूर्ति होते की नाही याचा आधार घेऊन परिभाषित केल्या
जात नाहीत; त्याऐवजी विशिष्ट सूत्रांची !!{derivability} किंवा
!!{nonderivability} यांचा आधार घेतला जातो. या दोन प्रकारच्या संकल्पना
एकमेकांशी जुळतात, हा एक महत्त्वाचा शोध होता. त्या जुळतात, हाच
\emph{निर्दोषता प्रमेय} आणि \emph{संपूर्णता प्रमेय} यांचा आशय आहे.
\end{explain}

\begin{defn}[!!^{derivability}]
!!^a{formula}~$!A$ हे $\Gamma$ \emph{पासून !!{derivable}} असते, म्हणजे
$\Gamma \Proves !A$, जर~$\Gamma$ पासून सुरू होऊन~$!A$ वर संपणारी
!!a{derivation} असेल.
\end{defn}

\begin{defn}[प्रमेये]
रिक्त संचापासून $!A$ ची !!a{derivation} असेल, तर !!^a{formula}~$!A$ हे
\emph{प्रमेय} असते. $!A$ हे प्रमेय असेल तर आपण $\Proves !A$ लिहितो आणि
ते प्रमेय नसेल तर $\Proves/ !A$ लिहितो.
\end{defn}

\begin{defn}[सुसंगतता]
!!{formula}s चा संच~$\Gamma$ हा $\Gamma\Proves/ \lfalse$ असेल तेव्हा आणि
केवळ तेव्हाच \emph{सुसंगत} असतो; अन्यथा तो \emph{विसंगत} असतो.
\end{defn}

\begin{prop}[परावर्तिता]
\ollabel{prop:reflexivity}
$!A \in \Gamma$ असेल, तर $\Gamma \Proves !A$.
\end{prop}

\begin{proof}
  एकटे !!{formula}~$!A$ ही~$\Gamma$ पासून $!A$ ची !!a{derivation} आहे.
\end{proof}

\begin{prop}[एकस्वनिकता]
\ollabel{prop:monotonicity}
$\Gamma \subseteq \Delta$ आणि $\Gamma \Proves !A$ असेल, तर $\Delta
\Proves !A$.
\end{prop}

\begin{proof}
$\Gamma$ पासून $!A$ ची कोणतीही !!{derivation} ही~$\Delta$ पासून $!A$ चीही
!!a{derivation} असते.
\end{proof}

\begin{prop}[संक्रमकता]
\ollabel{prop:transitivity}
$\Gamma \Proves !A$ आणि $\{!A\} \cup \Delta \Proves
!B$ असेल, तर $\Gamma \cup \Delta \Proves !B$.
\end{prop}

\begin{proof}
  $\{!A\} \cup \Delta \Proves !B$ असे समजा. मग~$\{!A\} \cup
  \Delta$ पासून $!B_1$, \dots, $!B_l = !B$ ही !!a{derivation} असते.
  त्या !!{derivation} मधील काही पायऱ्या योग्य असतात, कारण एखादा नियम
  त्यांपूर्वीच्या~$!B_i = !A$ या ओळीचा निर्देश करतो. गृहीतकानुसार,
  $\Gamma$ पासून~$!A$ ची !!a{derivation} असते, म्हणजे !!a{derivation}
  $!A_1$, \dots, $!A_k = !A$ अशी असते की प्रत्येक $!A_i$ हे स्वयंसिद्धक,
  $\Gamma$ चा !!a{element}, किंवा निगमन नियमाने योग्य ठरलेले असते. आता
  पुढील अनुक्रम विचारात घ्या:
  \[
  !A_1, \dots, !A_k = !A, !B_1, \dots, !B_l = !B.
  \]
  ही $\Gamma \cup \Delta$ पासून~$!B$ ची योग्य !!{derivation} आहे, कारण
  प्रत्येक $B_i = !A$ ला आता~$!A_k = !A$ ला समर्थित करणाऱ्या त्याच नियमाने
  समर्थन मिळते.
\end{proof}

लक्षात घ्या की विशेषतः याचा अर्थ असा होतो: $\Gamma \Proves !A$ आणि $!A
\Proves !B$ असेल, तर $\Gamma \Proves !B$. तसेच $!A_1,
\dots, !A_n \Proves !B$ असेल आणि प्रत्येक~$i$ साठी $\Gamma \Proves !A_i$
असेल, तर $\Gamma \Proves !B$, हेसुद्धा यावरून मिळते.

\begin{prop}
\ollabel{prop:incons}
$\Gamma$ विसंगत असतो तेव्हा आणि केवळ तेव्हाच, जेव्हा प्रत्येक~$!A$ साठी
$\Gamma \Proves !A$.
\end{prop}

\begin{proof}
सराव.
\end{proof}

\tagprob{FOL}
\begin{prob}
\olref[fol][axd][ptn]{prop:incons} सिद्ध करा.
\end{prob}
\tagendprob

\tagprob{notFOL}
\begin{prob}
\olref[pl][axd][ptn]{prop:incons} सिद्ध करा.
\end{prob}
\tagendprob

\begin{prop}[संहतता]
\ollabel{prop:proves-compact}
  \begin{enumerate}
  \item $\Gamma \Proves !A$ असेल, तर एखादा सांत उपसंच $\Gamma_0
    \subseteq \Gamma$ असा असतो की $\Gamma_0 \Proves !A$.
  \item $\Gamma$ चा प्रत्येक सांत उपसंच सुसंगत असेल, तर $\Gamma$ सुसंगत
    असतो.
  \end{enumerate}
\end{prop}

\begin{proof}
  \begin{enumerate}
    \item $\Gamma \Proves !A$ असेल, तर !!{formula}s चा सांत अनुक्रम
      $!A_1$, \dots,~$!A_n$ असा असतो की $!A \ident !A_n$ आणि प्रत्येक
      $!A_i$ हे एकतर तर्कशास्त्रीय स्वयंसिद्धक, $\Gamma$ चा !!a{element},
      किंवा आधीच्या !!{formula}s पासून विध्यात्मक अनुमानाने निष्पन्न झालेले
      असते. $\Gamma$ मध्ये असलेल्या त्या $!A_i$ चा संच~$\Gamma_0$ म्हणून
      घ्या. मग ही !!{derivation} त्याचप्रमाणे~$\Gamma_0$ पासूनची
      !!a{derivation} असते, आणि म्हणून $\Gamma_0 \Proves !A$.
    \item (1) मध्ये $!A
      \ident \lfalse$ हे विशेष प्रकरण घेतल्यावर मिळणारे हे निषेधात्मक उलट
      विधान आहे.
\end{enumerate}
\end{proof}

\end{document}
